\subsection{Saving}
\label{sec:saving}

Performance testing of update processing (see Chapter \ref{chap:results}) showed
that \snowman{} version 0.4 was very slow compared to other software. The reason
for this was examined, and found to be a large amount of database calls spread
out in the code. \snowman{} 0.4 saved every rule individually after a change was
found and the objects were updated. This principle is demonstrated in Listing
\ref{lst:wrongsave}, where the method tries to fetch, alter and save a
classification object, or create a new one if not found. This turned out to be a
bottleneck, as updates could take hours on low-end hardware. To solve this, we
split the methods of \pythoninline{Update} into two new classes:
\pythoninline{Parser} and \pythoninline{Updater}, where the first would only
parse the incoming data and send this to the latter which gathered all objects
and wrote these in bulk when the parsing was complete. As can be seen in Chapter
\ref{chap:results}, this had a significant impact on performance.

\begin{python}[caption=Updating a classification in \snowman{}
0.4,label=lst:wrongsave]
from core.models import RuleClass

def updateClassification(type, description, priority):
  try:
    # Update existing classification
    ruleclass = RuleClass.objects.get(classtype=type)
    ruleclass.description = description
    ruleclass.priority = priority
    ruleclass.save()
  except RuleClass.DoesNotExist:
    # Add new classification
    ruleclass = RuleClass.objects.create(classtype=type, ...)
\end{python}
As some of the objects refer to other objects in the
database through foreign-keys, it is vital to store the objects depending on
other objects after the objects they depend on are stored. A safe order to save
them has been identified and the objects are saved in the following order:
\begin{enumerate}
  \item Generators
  \item Classifications
  \item Reference-types
  \item RuleSets
  \item Rules
  \item RuleReferences
  \item Suppresses
  \item Filters
\end{enumerate}

For every item in the list above the procedure is approximately the same.
First, the objects received are attempted to be extracted from the database. The
response from the database lets us identify which objects are new and which are
changed. The changed objects can be updated and saved individually before all
the new objects are constructed and everything finally gets inserted into the
database in one bulk operation.
